Todas las entidades del SCDEE heredan de una jerarquía de tres clases
abstractas que aporta una base homogénea:

\begin{itemize}
  \item \texttt{UUIDModel} define un identificador primario \acs{uuid}
  v4. Frente al uso de enteros autoincrementales, los \acsp{uuid} evitan
  filtraciones de información (un identificador secuencial revela el
  número total de registros y el orden temporal de creación), son
  globalmente únicos a través de instancias y bases de datos, y son
  seguros frente a colisiones en escenarios de escalado horizontal.

  \item \texttt{TimestampedModel} extiende lo anterior con dos campos
  automáticos: \texttt{created\_at} (fijado en el \texttt{INSERT} y
  nunca modificado) y \texttt{updated\_at} (refrescado en cada
  \texttt{UPDATE}). Estos campos proporcionan una traza de auditoría
  implícita complementaria al log explícito
  (\cref{rf:security:audit-log}).

  \item \texttt{OrganizationOwnedModel} añade la clave foránea
  \texttt{organization} y dos \textit{managers}: \texttt{objects}
  (filtrado por organización mediante \texttt{TenantManager}) y
  \texttt{unfiltered} (sin filtro, para uso del superadministrador,
  tareas asíncronas y migraciones). La gran mayoría de modelos de
  negocio del SCDEE hereda de esta clase.
\end{itemize}

Las únicas excepciones a la herencia de \texttt{OrganizationOwnedModel}
son las entidades transversales del sistema: la propia
\texttt{Organization}, que es el contenedor multi-tenant, y
\texttt{AuditLog}, que cubre eventos cross-organización del
superadministrador.
